\section{Industrial Data Landscape: Lessons from the MVTec AD Benchmark}
\label{sec:dataset_landscape}

To validate the platform under authentic factory conditions, all algorithmic development and empirical benchmarking in this report are conducted on the MVTec Anomaly Detection (MVTec AD) benchmark, the established international standard for industrial computer vision research. The benchmark comprises 15 distinct manufacturing categories acquired in active production settings, dividing naturally into two structurally disparate inspection paradigms that together represent the overwhelming majority of shop-floor quality assurance tasks.

The first paradigm encompasses five continuous texture categories -- \emph{Carpet}, \emph{Grid}, \emph{Leather}, \emph{Tile}, and \emph{Wood} -- wherein the inspected surface constitutes a stationary, stochastic, or periodic material pattern. Defects in these categories manifest as localised breaks in spatial regularity: cut pile fibres, broken warp threads, surface scratches, or structural grain voids. The second paradigm encompasses ten rigid structural object categories -- \emph{Bottle}, \emph{Cable}, \emph{Capsule}, \emph{Hazelnut}, \emph{Metal Nut}, \emph{Pill}, \emph{Screw}, \emph{Toothbrush}, \emph{Transistor}, and \emph{Zipper}. These represent discrete assembled components where geometric contour, fixture pose, component count, and spatial assembly tolerances dictate conformity. Anomalous states here include structural fractures, missing sub-assemblies, mechanical deformations, and chemical surface discolourations.

\subsection{The Accuracy Paradox and Its Implications for Metric Selection}

The statistical reality of the MVTec AD benchmark exposes the fatal vulnerability of conventional classification metrics in industrial inspection. Across the benchmark, the median anomalous region accounts for merely $1.54\%$ of the total image area. This extreme spatial sparsity induces a severe mathematical pathology: a trivial majority-class classifier that marks every pixel as conforming achieves over $98\%$ overall pixel accuracy whilst failing to detect a single defect. In production metrology, this phenomenon represents the classic \emph{accuracy paradox}: aggregate accuracy provides zero indication of defect containment capability. Figure~\ref{fig:imbalance_eda} illustrates the distribution of defect region sizes and the acute class imbalance across all 15 benchmark categories.

\begin{figure}[tp]
    \centering
    \placeholder{2.0cm}{Bar chart or violin plot showing the distribution of anomalous pixel fractions across all 15 MVTec AD categories (x-axis: category name; y-axis: defect area as percentage of total image area). Each bar or violin should reflect the per-image defect area distribution for that category's test set. A horizontal dashed reference line at $1.54\%$ marks the dataset-wide median. The chart should make visually clear that the majority of categories have median defect areas well below $5\%$, motivating the rejection of pixel accuracy and pixel AUROC as primary metrics. Suggested palette: muted categorical colours per category group (textures vs.~objects).}
    \caption{Distribution of defect region magnitudes across the 15 MVTec AD categories, expressed as the anomalous pixel fraction of total image surface. The benchmark median of $1.54\%$ demonstrates the acute spatial imbalance: a trivial nominal classifier achieves $>98\%$ pixel accuracy whilst remaining completely blind to defects. This necessitates the exclusive adoption of PR-AUC and AUPIMO as primary evaluation metrics throughout this report.}
    \label{fig:imbalance_eda}
\end{figure}

The identical mathematical distortion compromises standard pixel-level AUROC. Because the Receiver Operating Characteristic integrates across the entire specificity domain, it is overwhelmingly dominated by the massive pool of true-negative conforming pixels. Consequently, any baseline system possessing rudimentary background segmentation can register a deceptive pixel AUROC of $>0.90$ whilst remaining virtually incapable of resolving micro-fractures occupying less than two per cent of the component surface. For this reason, the platform strictly rejects pixel AUROC as a primary commercial criterion. Instead, we mandate the Precision-Recall Area Under Curve (PR-AUC) and the Area Under the Per-Image Overlap curve (AUPIMO), integrated strictly over the ultra-low false-positive rate window $\text{FPR} \in [10^{-5}, 10^{-4}]$. These metrics discard irrelevant background mass, evaluating the system exclusively where commercial risk resides: on localised defect morphology.

\subsection{Heterogeneous Preprocessing Requirements Across Paradigms}

The fundamental dichotomy between stochastic textures and discrete structural objects dictates differentiated data preparation strategies. Continuous textures exhibit spatial translation invariance: a textile weave defect or grain irregularity can manifest at arbitrary spatial coordinates, and the model must exhibit uniform sensitivity across the entire field of view. Consequently, augmentation and preprocessing for texture categories prioritise spatial and affine transformations -- random cropping, controlled rotation, and elastic warping -- ensuring that the learned normal representation models the stationary stochastic process rather than spurious global position biases.

Conversely, discrete structural objects present the inverse challenge. Mechanical components are positioned within rigid mechanical jigs or conveyor nests, maintaining a largely deterministic spatial pose. However, their photometric characteristics fluctuate continuously across multi-shift operation: ambient factory illumination drifts, shadow angles shift due to conveyor vibrations, and camera sensor gain fluctuates with thermal variations. For structural objects, data preprocessing prioritises photometric normalisation -- contrast-limited adaptive histogram equalisation (CLAHE), gamma correction, and calibrated luminance jitter. This ensures feature embeddings remain invariant to environmental lighting noise whilst preserving the precise geometric alignment required for assembly tolerance verification.

\subsection{The Cold-Start Dilemma: Overcoming the Defect Collection Trap}

A decisive operational insight derived from the benchmark's composition is that supervised deep learning architectures are economically unfeasible in factory settings. The MVTec AD training partition consists exclusively of defect-free, nominal images; defective samples appear solely in the held-out evaluation split. This matches operational factory reality, where well-engineered production processes yield conforming parts in excess of $99.9\%$ of total volume. Defective components occur sporadically, exhibiting unpredictable morphologies that cannot be anticipated \emph{a priori}.

For manufacturing plant directors, deploying supervised vision models creates a prohibitive ``cold-start'' barrier. Assembling a statistically representative defect catalogue for supervised training would require a facility to intentionally fabricate, sort, and manually annotate tens of thousands of scrap parts over months of interrupted production before an inspection gate could be commissioned.

The platform eliminates this barrier entirely through self-supervised representation learning. By fitting the nominal manifold exclusively on standard production output, the platform establishes an empirical baseline of the conforming state (\emph{Soll-Zustand}). Any physical deviation from this distribution raises the continuous anomaly score. From an executive perspective, this architecture collapses the commissioning timeline: a newly introduced production line or tooling variant can be commissioned on day one using standard production parts, completely bypassing the costly defect collection trap.
